\section{Discussion}
\label{sec:discussion}

\subsection{What the reported numbers measure}
The three protocols differ in one respect, and that respect accounts for roughly
seventy per cent of the signal above chance. A practitioner reading $0.796$ and
a practitioner reading $0.585$ are reading the same model on the same corpus;
only the question being asked has changed. The first asks whether the detector
recognises functions resembling those it was trained on, which is a
well-specified question and not the one deployment poses. The second asks
whether it generalises to repositories it has never seen, which is.

This is not an argument that published numbers are wrongly computed. It is an
argument of the kind made more broadly about machine-learning-based
science~\citep{kapoor2023leakage}: the number alone does not identify which
question was answered, and the difference is large enough to change
conclusions. A benchmark
result reported without a statement of what the test partition was permitted to
share with training is therefore incomplete in a way that a reader cannot detect
or correct. Neither the mechanism nor the remedy is new to software engineering:
near-duplicate code is pervasive on the platforms these corpora are mined
from~\citep{lopes2017dejavu} and inflates model-of-code results when it crosses
the split~\citep{allamanis2019adverse}, and defect prediction established two
decades ago that cross-project evaluation asks a different and harder question
than within-project evaluation~\citep{zimmermann2009cross}. What our measurements
add is the size of the effect for this task and its decomposition into
duplication and repository overlap.

\subsection{Encoding choices are unreported and consequential}
How much a code property graph is measured to contribute is not a property of
the representation alone. It is a joint property of the representation and an
encoding decision that this literature does not vary and does not report.
Holding everything else fixed, restoring the reverse direction of each edge
moves the same graphs from undetectable to $+0.032$ ROC-AUC under
repository-disjoint evaluation, at no cost in model parameters. A second choice
that looks equally consequential, distinguishing the four relation types,
gains $0.002$ against a matched baseline (95\% CI $-0.005$ to $+0.010$).

The second result is worth stating precisely because the encoding it concerns
looks like the more meaningful one. That an edge carries data dependence rather
than syntactic containment is the kind of information a reader expects a code
property graph to supply, and it makes no measurable difference to any of the
three architectures able to read it.

The direction result is the more striking because the cause is mechanical rather
than statistical. Joern emits a directed graph; a network consuming it unaltered
lets each node aggregate only from its predecessors, which leaves the average
node able to reach under a fifth of its own function within three layers. The
fix is one line, and it is not applied because the encoding is treated as a
property of the input rather than as a design decision.

We note the asymmetry this creates in how the field allocates effort. A new
aggregation scheme is a publishable contribution; the orientation of the edges it
aggregates over is a detail of the data loader. On this corpus the second is
worth more than several of the first, and it is worth more than the relation
labels that motivate the representation in the first place. The same asymmetry has been reported for
graph classification generally, where architectural gains narrow substantially
once the surrounding protocol is held fixed~\citep{errica2020fair}.

\subsection{Aggregate metrics and the capability of interest}
The detectors examined here separate vulnerable from safe functions in aggregate
while failing, at chance, to separate a vulnerable function from its own patched
version. Those are different capabilities, and only the first is measured by the
metric the field reports.

The distinction matters for how results are read. A detector that recognises
which functions resemble known-vulnerable code is useful for triage and
prioritisation. A detector that identifies what specifically is wrong with a
function is what the framing of this literature, and its use of terms like
localisation and explanation, implies. Our measurements indicate that
improvements in the first do not imply progress in the second, and that
aggregate ROC-AUC cannot distinguish them. Reporting a paired evaluation
alongside it costs one additional column and closes that gap.

\subsection{Implications for practice}
Three recommendations follow directly from the measurements. Evaluations should
ship an audit of what the split separates, since the cost of omitting one is
both large and currently invisible. Ablations should report seed counts and a test of significance, because
run-to-run variation here is of the same magnitude as the effects being
measured and the appropriate procedures are
well established~\citep{arcuri2011practical}. And studies reporting a null result for a component should include a
positive control establishing that the model can learn a label defined by that
component, without which a null cannot be distinguished from an instrument that
never read it.

A fourth recommendation concerns the training procedure itself. Whether a study
trains a single model or an ensemble, and how it selects the reported epoch, is
worth more than several of the effects such studies report, and is rarely
stated. We found that selecting the epoch on data sharing repositories with the
training set costs ROC-AUC under a repository-disjoint protocol, which is the
same leakage the split is designed to prevent, reappearing inside the training
loop. That model validation technique is itself a first-order determinant of
reported performance is a finding defect prediction reached
earlier~\citep{tantithamthavorn2017empirical}.
